\documentclass[11pt]{article}
\usepackage[margin=1in]{geometry}
\usepackage{amsmath,amssymb,amsthm}
\usepackage{enumitem}
\usepackage{graphicx}
\usepackage[hidelinks]{hyperref}

\newtheorem{theorem}{Theorem}
\newtheorem{lemma}{Lemma}

\newcommand{\R}{\mathbb R}
\newcommand{\lp}{\ell_{20}^{2}}
\newcommand{\norm}[1]{\left\|#1\right\|}
\newcommand{\la}{\lambda}

\title{A Counterexample to the \(L_p\) Projection-Constant Equality Conjecture}
\author{Run \texttt{fa\_banach\_001}, agent \texttt{agent\_lane\_16}}
\date{2026-06-28}

\begin{document}
\sloppy
\maketitle

\section*{Source Problem}

The source is G. M. Ivanov, \emph{Modulus of supporting convexity and
supporting smoothness}, arXiv:1503.08912. In Section 6 the paper defines
\[
\xi_X=
\sup_{\norm{x}=1,\ \norm{y}=1}\sup_{p\in J_1(y)}
\norm{x-\langle p,x\rangle y},
\]
the maximal Lipschitz constant of the metric projection onto a hyperplane.
Theorem 4 proves
\[
\xi_X\leq
\frac{1}{1-\lambda_X^+\!\left((1-\lambda_X^-(1))/2\right)}.
\]
Conjecture 1 states that this right-hand side becomes an equality for
\(L_p\), \(p\in(1,\infty)\).

\begin{center}
\includegraphics[width=0.9\linewidth]{figures/open_problem_crop.png}
\end{center}

\section*{Proof Intuition}

The conjecture is already false for the two-atom \(L_p\) space
\(X=\ell_{20}^2\). In two dimensions, every support projection has rank one
and its norm can be written explicitly. This gives a global analytic bound
\(\xi_X\leq 2^{9/10}\).

The source upper expression is larger. To see this, first make
\(\lambda_X^-(1)\) almost zero using a very asymmetric tangent pair. This
forces the argument \((1-\lambda_X^-(1))/2\) of \(\lambda_X^+\) to be almost
\(1/2\). At radius \(0.4999995\), the symmetric tangent pair in \(\ell_{20}^2\)
already yields a \(\lambda^+\)-value whose reciprocal upper expression exceeds
\(2^{9/10}\). Therefore the source upper estimate is strict.

\section*{Counterexample}

\begin{theorem}
Conjecture 1 of Ivanov's paper is false. More precisely, for
\(X=\ell_{20}^{2}\),
\[
\xi_X<
\frac{1}{1-\lambda_X^+\!\left((1-\lambda_X^-(1))/2\right)}.
\]
\end{theorem}

\begin{proof}
Write \(q=20/19\). Let \(y=(a,b)\in S_X\), so
\(|a|^{20}+|b|^{20}=1\), and let \(j(y)=(\operatorname{sgn}(a)|a|^{19},
\operatorname{sgn}(b)|b|^{19})\) be the norming functional of \(y\).
By signs and symmetry it is enough to display the computation for
\(a,b\geq0\).

The projection \(P_yx=x-j(y)(x)y\) has range \(\ker j(y)\), which is spanned
by \(u=(b^{19},-a^{19})\). Since
\(\det(u,y)=a^{20}+b^{20}=1\), we have
\[
P_y(s,t)=(bs-at)u .
\]
Consequently
\[
\norm{P_y}
=\left(a^{380}+b^{380}\right)^{1/20}
 \left(a^{q}+b^{q}\right)^{1/q}.
\]
The first factor is at most \(1\). For the second factor, Holder's inequality
or the standard comparison of finite-dimensional \(\ell_r\)-norms gives
\[
\left(a^q+b^q\right)^{1/q}\leq
2^{1/q-1/20}\left(a^{20}+b^{20}\right)^{1/20}
=2^{18/20}=2^{9/10}.
\]
Taking the supremum over \(y\) gives
\begin{equation}\label{eq:xi-upper}
\xi_X\leq 2^{9/10}.
\end{equation}

We next show that the right-hand side of Ivanov's estimate is strictly larger
than \(2^{9/10}\). Put \(t=9/20\), \(\varepsilon=10^{-6}\),
\[
A=(1+t^{20})^{1/20},\qquad
B=(1+t^{380})^{1/20},
\]
and
\[
x=\frac{(1,t)}{A},\qquad y=\frac{(t^{19},-1)}{B}.
\]
Then \(x,y\in S_X\), and \(y\) is quasi-orthogonal to \(x\), because the
norming functional of \(x\) is \(A^{-19}(1,t^{19})\), which annihilates \(y\).
A direct estimate gives
\[
\left|
\frac{1-\varepsilon}{A}+\frac{t^{19}}{B}
\right|^{20}
+
\left|
\frac{(1-\varepsilon)t}{A}-\frac{1}{B}
\right|^{20}
<1 .
\]
The verifier in \texttt{code/check\_counterexample.py} evaluates the left
side as \(0.9999914532\ldots\). The same script evaluates
\(\norm{x+y}_{20}^{20}\) as \(1.0000114530\ldots\). Hence the point
\((1-\varepsilon)x+y\) lies inside the unit ball, while \(x+y\) lies outside
it. By continuity of the line section,
\[
\lambda_X^-(1)<\varepsilon,
\qquad
r:=\frac{1-\lambda_X^-(1)}2>\frac{1-\varepsilon}{2}=:r_0.
\]

Now take the symmetric quasi-orthogonal pair
\[
x_0=2^{-1/20}(1,1),\qquad y_0=2^{-1/20}(1,-1).
\]
For this pair, writing \(u=1-\lambda\), the equation
\(\norm{(1-\lambda)x_0+r_0y_0}_{20}=1\) is
\[
(u+r_0)^{20}+|u-r_0|^{20}=2.
\]
Let \(u_0=2^{-9/10}\). A direct strict check gives
\[
(u_0+r_0)^{20}+|u_0-r_0|^{20}>2
\]
(the verifier gives \(2.0241430960\ldots\)). Since the left-hand side is
strictly increasing for \(u>r_0\), the corresponding root satisfies
\(u<u_0\), i.e. the symmetric pair gives
\[
\lambda_X^+(r_0)>1-u_0.
\]
As \(\lambda_X^+\) is increasing and \(r>r_0\),
\[
\lambda_X^+(r)>1-2^{-9/10}.
\]
Therefore
\[
\frac{1}{1-\lambda_X^+(r)}>2^{9/10}\geq \xi_X
\]
by \eqref{eq:xi-upper}. This is a strict failure of the equality conjectured
for \(L_p\).
\end{proof}

\section*{Verification Notes}

The proof is analytic except for two scalar strict inequalities involving
fixed decimal/rational parameters. They are also easy to check by hand with
interval arithmetic; the included Python script records a reproducible
high-precision check:
\[
1.0000114530\ldots>1,\qquad
0.9999914532\ldots<1,\qquad
2.0241430960\ldots>2.
\]

Bounded novelty check: the run indexes were searched for arXiv:1503.08912,
supporting convexity/smoothness, hyperplane projection constants, and
Birkhoff-James terminology. Web searches on 2026-06-28 for the exact title,
the \(L_p\) conjecture phrase, and hyperplane projection constants found the
source arXiv record and the related later Ivanov--Martini paper
arXiv:1609.01587, but no explicit later resolution of this conjecture.
The run already contains a counterexample packet for a different conjecture in
arXiv:1609.01587.

\section*{Human Review Recommendation}

Review as a likely valid counterexample. The main checks are the rank-one
projection formula in \(\ell_{20}^{2}\), the orientation of
quasi-orthogonality relative to the source definitions, and the two strict
scalar inequalities above.

\begin{thebibliography}{9}

\bibitem{Ivanov}
G. M. Ivanov,
\emph{Modulus of supporting convexity and supporting smoothness},
arXiv:1503.08912, 2015.

\bibitem{IvanovMartini}
G. Ivanov and H. Martini,
\emph{New Moduli for Banach Spaces},
arXiv:1609.01587, 2016.

\end{thebibliography}

\end{document}
